% !PRINTED 296
\weberpage{296}
%<<CONT
on the idea that underlies the concept of the characteristic\,\footnote{\gesperrt{Gauss}, Demonstratio nova theorematis omnem functionem algebraicam rationalem integram unius variabilis in factores reales primi vel
secundi gradus resolvi posse (1799). Beiträge zur Theorie der algebraischen
Gleichungen (1849). \gesperrt{Kronecker}, Monatsberichte der Berliner Akademie,
21. Februar 1878.}. We therefore fittingly close this section with an
exposition of this proof.

Let, then, $z = x + yi$ and
\weq{1}{F(z) = \varphi(x, y) + i\,\psi(x, y)}
be an integral rational function of $z$ of degree $n$ with real
or complex coefficients, in which $z^n$ has the coefficient 1,
thus
\weq{2}{F(z) = z^n + a_1 z^{n-1} + a_2 z^{n-2} + \cdots + a_n.}

We assume again that $F(z)$ and $F'(z)$ have no common
divisor, and the proof that $F(z)$ vanishes for $n$
values of $z$ now rests on the fact that, if we
choose for $f$ a circle of sufficiently large radius, the
characteristic of the system of functions $f, \varphi, \psi$ can be determined
directly and is found to be equal to $n$. It then follows by \S.\,96
that in the circle so chosen there lie $n$ points $z$ at which
$F(z)$ vanishes.

We make use of polar coordinates, and set, since
$a_1, a_2, a_3 \dots a_n$ may also be complex,
\[
a_1 = p_1 e^{iq_1},\quad a_2 = p_2 e^{iq_2}, \dots a_n = p_n e^{iq_n}
\]
\[
z = R e^{i\vartheta},
\]
in which $p_1, p_2, \dots p_n, R$ are positive and the angles $q_1, q_2, \dots q_n, \vartheta$
can be chosen in any interval of extent $2\pi$.
Hereby we obtain
\weq{3}{
\begin{aligned}
\varphi(x, y) &= R^n \cos n\vartheta + p_1 R^{n-1} \cos[(n-1)\vartheta + q_1] \\
&\quad + \cdots + p_n \cos q_n \\[4pt]
\psi(x, y) &= R^n \sin n\vartheta + p_1 R^{n-1} \sin[(n-1)\vartheta + q_1] \\
&\quad + \cdots + p_n \sin q_n,
\end{aligned}
}
and we also form the derivatives taken with respect to $\vartheta$
%CONT>>
